% !PRINTED 145
\weberpage{145}
%<<CONT
then the coefficients $S_0,S_1\ldots S_\mu$ are whole \gesperrt{symmetric} functions of the $n-1$ variables $\alpha_2,\ldots\alpha_n$.

If we denote the fundamental symmetric functions of these latter variables by $a'_1,a'_2\ldots a'_{n-1}$, then, by our supposition, we can express the coefficients $S_0,S_1\ldots S_\mu$ rationally in terms of these. But by \S\,42, (2) and (3),
\[
\begin{array}{l}
 a'_1=f_1(\alpha_1)=\alpha_1+a_1\\
 a'_2=f_2(\alpha_1)=\alpha_1^2+a_1\alpha_1+a_2\\
 a'_3=f_3(\alpha_1)=\alpha_1^3+a_1\alpha_1^2+a_2\alpha_1+a_3\\
 \makebox[23em][l]{\dotfill},
\end{array}
\]
that is, the coefficients $S_0,S_1\ldots S_\mu$ can be expressed wholly and rationally in terms of $\alpha_1,a_1,a_2\ldots a_n$. If, therefore, after introducing these expressions, we arrange (2) anew according to powers of $\alpha_1$, we obtain
\weq{3}{S=A_0\alpha_1^m+A_1\alpha_1^{m-1}+\cdots+A_{m-1}\alpha_1+A_m,}
where in general $m$ will be an exponent different from $\mu$, and the coefficients $A_0,A_1\ldots A_m$ depend wholly and rationally on $a_1,a_2\ldots a_n$. We put again
\weq{4}{\Phi(x)=A_0x^m+A_1x^{m-1}+\cdots+A_{m-1}x+A_m}
and divide $\Phi(x)$ by
\[
\begin{aligned}
 f(x)&=x^n+a_1x^{n-1}+a_2x^{n-2}+\cdots+a_n\\
     &=(x-\alpha_1)(x-\alpha_2)\ldots(x-\alpha_n).
\end{aligned}
\]
This gives a quotient and a remainder which, with respect to $x$, is at most of degree $n-1$. We therefore put
\weq{5}{\Phi(x)=Qf(x)+\psi(x),}
\weq{6}{\psi(x)=C_0x^{n-1}+C_1x^{n-2}+\cdots+C_{n-2}x+C_{n-1},}
and here $C_0,C_1\ldots C_{n-1}$ are whole rational functions of $a_1,a_2\ldots a_n$ in which, if $S$ in its original form contains no fractional coefficients, no fractions occur either.

But now, since $f(\alpha_1)$ vanishes and $S=\Phi(\alpha_1)$, it follows from (5) that
\weq{7}{S=\psi(\alpha_1)}
and here, since $S$ is symmetric, $\alpha_1$ may be replaced by $\alpha_2,\alpha_3\ldots\alpha_n$. This gives the following $n$ equations:
%CONT>>
